% Generated draft: requires mathematical and source review.
\documentclass[11pt]{article}
\usepackage[margin=25mm]{geometry}
\usepackage{fontspec}
\setmainfont{Latin Modern Roman}
\usepackage{amsmath,amssymb,mathtools,unicode-math}
\setmathfont{Latin Modern Math}
\usepackage{enumitem,needspace,xurl,hyperref}
\setlength{\parindent}{0pt}
\setlength{\parskip}{5pt}
\setlength{\emergencystretch}{3em}
\newcommand{\N}{\mathbb{N}}
\newcommand{\Z}{\mathbb{Z}}
\newcommand{\Q}{\mathbb{Q}}
\newcommand{\R}{\mathbb{R}}
\newcommand{\C}{\mathbb{C}}
\newcommand{\F}{\mathbb{F}}
\newcommand{\A}{\mathbb{A}}
\newcommand{\PP}{\mathbb{P}}
\newcommand{\E}{\mathbb{E}}
\newcommand{\eps}{\varepsilon}
\newcommand{\dd}{\,\mathrm{d}}
\begin{document}
\section*{UnsolvedMath ID: 30004227}
\textbf{Source problem code:} catalogue entry
\subsection{Definitions and mathematical statement}
The supplied record poses three related research questions, but does not specify enough conventions to determine a unique mathematical formulation.

For the topology question, let \(k\) denote the number of sites and \(d\) the dimension. The source asks for the topological structure of tropical bisectors of \(k\) points in every dimension \(d\). It does not define the ambient tropical space, its distance function, the meaning of a bisector of several sites, or any general-position assumptions. For example, if an ambient space \(X_d\) with distance \(\delta_d\) were specified and ``bisector'' meant the common equal-distance locus, its definition would be
\[
 B_{\delta_d}(p_1,\ldots,p_k)
 =\{x\in X_d:\delta_d(x,p_1)=\cdots=\delta_d(x,p_k)\},
 \qquad p_1,\ldots,p_k\in X_d.
\]
This is only a conditional definition: the supplied record does not establish that this convention is intended. In particular, it does not distinguish this locus from a union of pairwise bisectors or a locus where the minimum distance to the sites is attained more than once. Subject to resolving these choices, the request is to determine the topology for all admissible \(k\), all dimensions \(d\), and all configurations covered by the intended hypotheses. Neither the precise parameter ranges nor the intended type of topological classification is supplied.

For the algorithmic question, the source asks whether an expected construction running time
\[
 O(nd\log n)
\]
can be improved when possible. The meaning of \(n\), its relation to \(k\), the object constructed, the input and output representations, the computational model, the source of randomness, and any restrictions on the input are unspecified. Thus the displayed bound is retained as a supplied benchmark, without asserting a particular algorithmic theorem or a precise improvement target.

For the curvature question, the source asks for a useful notion of tropical curvature for arbitrary polygonal space curves. No ambient space, class of admissible polygonal curves, boundary conditions, curvature codomain, or required properties of such a notion are specified. The word ``arbitrary'' supplies no stated restriction to generic, embedded, closed, or otherwise special curves. The request therefore remains a definition problem whose intended requirements need clarification; no curvature formula is asserted here.
\subsection{Short English statement}
Determine the topology of tropical bisectors for all numbers of sites and dimensions, investigate improvements to the supplied expected \(O(nd\log n)\) construction bound, and develop a useful tropical curvature for arbitrary polygonal space curves. The supplied record leaves essential definitions and algorithmic hypotheses unspecified, so this is a qualified research-program statement requiring human mathematical review.
\subsection{Sources}
\begin{itemize}
\item Local UnsolvedMath registry, permanent ID \(30004227\), source problem code \texttt{catalogue entry}, title ``Topology Algorithms and Curvature of Tropical Bisectors,'' in the Oberwolfach Reports Open Problems collection. The supplied record lists proposed year 2019 and names no proposer.
\item Supplied source citation: \emph{Geometric, Algebraic, and Topological Combinatorics (2020)}, Oberwolfach Reports, pp.\ 2395--2472, DOI \texttt{10.4171/owr/2019/39}. \url{https://doi.org/10.4171/owr/2019/39}. The record describes Oberwolfach Reports as documenting workshops and their open-problem discussions. The supplied citation uses 2020, whereas the DOI and proposed-year field refer to 2019; this dating has not been independently checked.
\item Supplied literature source: Criado, Joswig, and Santos, \url{https://doi.org/10.1007/s10208-021-09538-4}. The record supplies no title, precise result references, or individual problem attribution for this citation.
\item Literature assessment dated exactly \texttt{2026-08-22}: the supplied curation labels the problem open and reports results for \(k=2\) and \(k=3\), examples for \(k=4\), and remaining open topology, algorithmic, and curvature questions for general parameters. Its listed sources are the Oberwolfach report and Criado, Joswig, and Santos, linked above. This is supplied dated curation derived from the OWR material, not an independently verified literature assessment; the record explicitly requests independent verification of time-sensitive status and cited claims.
\item The registry's supplied statement-status label is \texttt{corrected\_verified}; its accompanying explanation concerns preservation of mathematical content during notation, typography, and sentence-structure repairs. This label is recorded as source metadata and does not establish mathematical verification of this draft.
\end{itemize}
\end{document}
